\documentclass[report]{itvest}

% Document-specific metadata
\title{Portfolio 3: Observability with Knowledge Graphs}
\reportsubtitle{Case: Short Anonymised Case Name}
\author{Student Name} % group: \author{Name 1 \and Name 2 \and Name 3}
\date{\today}

% Portfolios use sections only, numbered 1, 2, ...
\renewcommand{\thesection}{\arabic{section}}

\begin{document}

\maketitle

\begin{checkit}
  \textbf{Assessment criteria} (delete this box before submitting):
  \begin{itemize}
    \item The graph covers both the system and lineage part and the domain
      part.
    \item The queries answer concrete observability questions for the case.
    \item Governance rules are concrete and consistent with the P2
      data-placement table.
    \item Drift and quality signals are defined, measurable, linked to the
      graph and tied to an action.
    \item At least one feedback loop reconfigures the running system,
      designed concretely enough to be built.
  \end{itemize}
\end{checkit}

\section{Data Flow}

Show how data flows from its sources through the pipelines into the
knowledge graph and the model, with the tier boundaries from P2.

\begin{figure}[h]
  \centering
  % Replace the box with: \includegraphics[width=\textwidth]{data-flow}
  \fbox{\parbox[c][7cm][c]{0.9\textwidth}{\centering
    Replace with the data-flow diagram, tier boundaries marked.}}
  \caption{Data flow.}
  \label{fig:data-flow}
\end{figure}

Estimate what data feeds the graph, at what volume and rate, and whether it
arrives in batches or as a stream~\cite{huyen_designing_2022}. Show the
calculation, e.g.
\[
  \text{volume per day} =
  \text{sources} \times \text{records per second} \times
  \text{bytes per record} \times 86\,400 .
\]

\begin{table}[h]
  \centering
  \begin{tabularx}{\textwidth}{Xllll}
    \toprule
    Data source & Batch / stream & Rate & Size per record & Volume per day \\
    \midrule
    \ldots & \ldots & \ldots & \ldots & \ldots \\
    \ldots & \ldots & \ldots & \ldots & \ldots \\
    \bottomrule
  \end{tabularx}
  \caption{Volume and rate estimate.}
  \label{tab:volume}
\end{table}

\section{Knowledge-Graph Model}

Model your system and its domain as a knowledge graph. The
\emph{system and lineage} part describes components, data sources,
datasets, features, model versions and deployments, and how a deployed
model can be traced back to the data it was trained on. The \emph{domain}
part gives the observations meaning, e.g.\ assets, sensors, the properties
they observe, and the observations themselves.

\begin{figure}[h]
  \centering
  % Replace the box with: \includegraphics[width=\textwidth]{graph-schema}
  \fbox{\parbox[c][7cm][c]{0.9\textwidth}{\centering
    Replace with the graph schema: classes as nodes, relations as edges.}}
  \caption{Knowledge-graph schema.}
  \label{fig:schema}
\end{figure}

\begin{table}[h]
  \centering
  \begin{tabularx}{\textwidth}{llX}
    \toprule
    Class or relation & Part & Meaning \\
    \midrule
    \ldots & System and lineage & \ldots \\
    \ldots & Domain & \ldots \\
    \bottomrule
  \end{tabularx}
  \caption{Classes and relations.}
  \label{tab:schema}
\end{table}

\subsection{ADR-4: Building, Storing and Querying the Graph}

\begin{description}
  \item[Status:] Proposed / accepted.
  \item[Context:] What the graph must answer, how often it changes, and the
    constraints from P2.
  \item[Alternatives:] The options you considered.
  \item[Decision:] The chosen alternative.
  \item[Consequences:] The trade-off that decided it, and what gets harder.
\end{description}

\section{Observability Questions}

Give at least three questions the graph must answer, e.g.\ ``what is
affected if this sensor drifts?'' or ``which data trained the model in
production?'', each with the query that answers it and an example result.
Put long queries in an appendix.

\begin{table}[h]
  \centering
  \begin{tabularx}{\textwidth}{XXX}
    \toprule
    Question & Query (summary) & Example result \\
    \midrule
    \ldots & \ldots & \ldots \\
    \ldots & \ldots & \ldots \\
    \ldots & \ldots & \ldots \\
    \bottomrule
  \end{tabularx}
  \caption{Observability questions.}
  \label{tab:questions}
\end{table}

\section{Data Governance}

State who owns which data, how long it is kept, and which regulations
apply (e.g.\ GDPR, NIS2, EU AI Act), and what they require of the design.
The locations must match the P2 data-placement table. The governance
information may also be kept in the graph.

\begin{table}[h]
  \centering
  \begin{tabularx}{\textwidth}{XlllX}
    \toprule
    Data item & Owner & Retention & Location & Applicable rule \\
    \midrule
    \ldots & \ldots & \ldots & \ldots & \ldots \\
    \ldots & \ldots & \ldots & \ldots & \ldots \\
    \bottomrule
  \end{tabularx}
  \caption{Data governance.}
  \label{tab:governance}
\end{table}

\section{Observability}

\subsection{Data Quality}

Describe the checks on incoming data (e.g.\ schema, ranges, missing
values) and where they run.

\subsection{Drift and Proxy Metrics}

Describe how drift is detected while labels arrive late, and which proxy
metrics stand in for model quality until they do.

\subsection{Service and Model Quality}

Separate service SLOs from model-quality metrics, and link each metric to
the graph nodes it concerns.

\begin{table}[h]
  \centering
  \begin{tabularx}{\textwidth}{XllXX}
    \toprule
    Metric & Kind & Target & Measured where & Graph node \\
    \midrule
    \ldots & Service SLO & \ldots & \ldots & \ldots \\
    \ldots & Model quality & \ldots & \ldots & \ldots \\
    \ldots & Data quality / drift & \ldots & \ldots & \ldots \\
    \bottomrule
  \end{tabularx}
  \caption{Observability metrics.}
  \label{tab:metrics}
\end{table}

\subsection{Proactive Feedback}

Tie each signal to an action: retrain, roll back or reconfigure. At least
one action must reconfigure the running system without a manual redeploy,
e.g.\ switching the model version, adjusting a threshold, or falling back
to a simpler model. Use the graph to show what the action affects.

\begin{table}[h]
  \centering
  \begin{tabularx}{\textwidth}{XXX}
    \toprule
    Signal & Threshold & Action \\
    \midrule
    \ldots & \ldots & \ldots \\
    \bottomrule
  \end{tabularx}
  \caption{Signals and actions.}
  \label{tab:feedback}
\end{table}

\section{Evidence}

State whether the evidence comes from your own case or from the reference
project, and show the populated graph and the queries running. Dashboards
and drift evidence from week 6 may be added in the final report.

\section{AI Usage}

Briefly describe your use of AI tools on this assignment:

\begin{itemize}
  \item \emph{Code generation:} \ldots
  \item \emph{Grammar and writing:} \ldots
  \item \emph{Clarification of technical aspects:} \ldots
\end{itemize}

\bibliographystyle{IEEEtran}
\bibliography{bib}

\end{document}
